\subsection*{Sprint 1: Configuración del repositorio}
\addcontentsline{toc}{subsection}{Sprint 1: Configuración del repositorio}

El primer sprint se enfoca en establecer los cimientos técnicos del proyecto
\textbf{Muncher} mediante la creación de un repositorio en GitHub, la
configuración de un proyecto base utilizando Vite y React, y la implementación
de controles de calidad automatizados a través de GitHub Actions con Prettier y
ESLint.

La incorporación de estas herramientas de calidad desde las etapas más
tempranas del desarrollo es fundamental para el éxito del proyecto. Como
establece el principio de que "la prevención es más efectiva que la
corrección", implementar controles de formato de código y análisis estático
cuando la base de código es mínima permite establecer estándares consistentes
desde el inicio. Esto reduce significativamente la deuda técnica futura, ya que
cuanto menos código fuente existe, menor es la superficie de potenciales
errores que podrían requerir corrección posteriormente, y más sencillo resulta
mantener la consistencia y calidad del código a medida que el proyecto crece.

\subsubsection*{Objetivos iniciales}

\begin{itemize}
	\item Configuración del repositorio en GitHub.
	\item Creación de un proyecto base utilizando Vite y React.
	\item Implementación de controles de calidad automatizados a través de GitHub Actions
	      con Prettier y ESLint.
\end{itemize}

\subsubsection*{Historias de usuario completadas}

\begin{center}
	\begin{tabularx}{\textwidth}{|l|X|c|}
		\hline
		\textbf{Identificador} & \textbf{Nombre}                               & \textbf{Puntos de historia} \\
		\hline
		\usref{US-CI-01}       & Creación del repositorio de código            & 1                           \\
		\hline
		\usref{US-CI-02}       & Creación de una SPA usando Vite y React       & 1                           \\
		\hline
		\usref{US-CI-03}       & Controles de calidad con ESLint y Prettier    & 3                           \\
		\hline
		\usref{US-CI-04}       & Corrección automática del código de front-end & 5                           \\
		\hline
	\end{tabularx}
\end{center}

\subsubsection*{Retrospectiva}

El sprint finalizó con éxito, cumpliendo todos los objetivos planteados. El
repositorio en GitHub quedó configurado con una base sólida mediante Vite y
React, junto con la integración de controles de calidad automatizados que
aseguran un código consistente desde el inicio del proyecto.

La primera dificultad técnica consistió en aplicar la corrección automática del
código en un entorno \textit{monorepo} que reúne tecnologías distintas, cada una
con sus propias herramientas y convenciones de estilo. El front-end en React se
formatea con Prettier y se analiza con ESLint, mientras que el back-end en
Python y las definiciones de las imágenes de contenedor exigen sus propios
formateadores y analizadores estáticos; someter todos los archivos a todas las
herramientas no tiene sentido, de modo que cada cambio debe dirigirse
únicamente a las que le corresponden.

Para resolverlo se adoptó \texttt{pre-commit} como gestor de los controles en la
raíz del repositorio.\footnote{Documentación oficial: \url{https://pre-commit.com/}}
Esta herramienta declara los controles en un único archivo de configuración y
asocia cada uno a un patrón de rutas, por lo que es capaz de encaminar los
archivos modificados hacia las herramientas pertinentes sin necesidad de lógica
propia: los controles generales —espacios sobrantes, saltos de línea finales,
validez de los archivos YAML y JSON o presencia de marcas de conflictos de
fusión— se aplican a todo el repositorio, mientras que los específicos de cada
lenguaje se limitan a los directorios que les conciernen.

La segunda dificultad, de mayor calado, fue la dependencia de las herramientas
instaladas en la máquina de cada desarrollador. Un control solo resulta útil si
se ejecuta siempre y del mismo modo, y ninguna de esas dos condiciones se
cumple cuando la cadena de herramientas se instala manualmente: basta con que
una de ellas falte en un equipo para que el control se omita precisamente allí
donde hacía falta. A ello se añade la divergencia de versiones, ya que cada
desarrollador y la propia integración continua pueden ejecutar versiones
distintas del mismo analizador, con lo que un cambio se aprueba en local y se
rechaza en el servidor. El problema se agrava en un \textit{monorepo}, donde no
existe una única cadena de herramientas que instalar, sino una por tecnología.

La solución consistió en trasladar al interior de contenedores la ejecución de
todo aquello que el desarrollo requiere, tanto los controles como la propia
aplicación, conforme al criterio expuesto en
\hyperref[sec:entorno-local]{Entorno local}. En el caso de los controles, una
única imagen reúne las cadenas de herramientas necesarias para aplicarlos, con
los \textit{runtimes} y las dependencias fijados en su definición. El
\textit{hook} de git se limita a delegar la ejecución en el contenedor, y la
integración continua construye y emplea esa misma imagen en lugar de instalar
herramientas por su cuenta, de modo que los resultados obtenidos en local y en
el servidor dejan de poder divergir.

La aplicación se trata del mismo modo, aunque con una diferencia importante.
Cada componente dispone de su propia imagen, que contiene su \textit{runtime} y
sus dependencias, pero no el código fuente: este permanece en la máquina del
desarrollador y se monta en el contenedor durante la ejecución. De esta manera
la imagen constituye una instantánea del entorno, mientras que el código sigue
editándose con las herramientas habituales y los cambios se reflejan de
inmediato, sin necesidad de reconstruir ni reiniciar nada. Los servicios que
componen el sistema se declaran en una única composición y comparten una misma
red, lo que les permite localizarse entre sí por su nombre y reproduce en local
la topología del entorno desplegado. El resultado es que el desarrollador solo
necesita git y Docker: ni los \textit{runtimes} de los
lenguajes empleados, ni los gestores de dependencias, ni las herramientas de
análisis se instalan en el sistema anfitrión.

Gracias a esta solución, el proceso de desarrollo se beneficia de una mayor
calidad y consistencia, reduciendo la necesidad de correcciones manuales,
evitando el riesgo de introducir errores por cambios de formato y garantizando
que los controles se comporten de igual modo en cualquier entorno.
